\documentclass{article}

\usepackage[margin=1in]{geometry}
\usepackage{amsmath,amsthm,amssymb}
\usepackage[spanish]{babel}
\usepackage{enumitem}

\theoremstyle{definition}
\newtheorem{definition}{Definición}[section]

% Number sets
\newcommand{\R}{\mathbb{R}}
\newcommand{\Z}{\mathbb{Z}}
\newcommand{\N}{\mathbb{N}}
\newcommand{\Q}{\mathbb{Q}}
\newcommand{\C}{\mathbb{C}}
\newcommand{\K}{\mathbb{K}}

% Functional analysis operators
\newcommand{\norm}[1]{\left\lVert #1 \right\rVert}       % Norm
\newcommand{\seminorm}[1]{\left\lVert #1 \right\rVert}   % Seminorm (same symbol)
\newcommand{\cl}[1]{\overline{#1}}                        % Closure
\newcommand{\Int}[1]{{#1}^\circ}                          % Interior
\newcommand{\bola}[2]{B(#1;#2)}                           % Ball centered at #1 with radius #2
\newcommand{\EN}[1][E]{(#1,\norm{\cdot})}                 % Normed space (E, ||.||); use \EN or \EN[F]

% Useful shortcuts
\newcommand{\sii}{\iff}                          % Si y sólo si
\newcommand{\imp}{\implies}                      % Implies
\newcommand{\setdiff}{\smallsetminus}            % Set difference
\newcommand{\diam}{\text{diam}}

% Exercise environment
\newenvironment{ejercicio}[1]{\begin{trivlist}
\item[\hskip \labelsep {\bfseries Ejercicio}\hskip \labelsep {\bfseries #1.}]}{\end{trivlist}}

\setlist[enumerate,1]{label=\alph*), leftmargin=*, itemsep=2pt}
\setlist[enumerate,2]{label=\roman*., leftmargin=*, itemsep=2pt}

\begin{document}

\title{Análisis Funcional: Práctica II}
\author{Bustos Jordi\\Funcionales y Operadores en espacios de Banach}

\maketitle

\begin{ejercicio}{1}
    Sea \(E\) un espacio vectorial. Demostrar que si \(\varphi:E\to\C\) es un funcional lineal no nulo entonces \(\operatorname{im}(\varphi)=\C\).
\end{ejercicio}
\begin{proof}

\end{proof}

\vspace{0.5cm}

\begin{ejercicio}{2}
    Sean \(E\) un espacio vectorial normado sobre \(\R\) y \(\varphi:E\to\R\) un funcional lineal. Probar
    \begin{enumerate}
        \item Si \(\varphi\) no es acotado entonces toma todos los valores reales en cualquier entorno del \(0\).
        \item Si \(\varphi\) no es acotado entonces \(\ker(\varphi)\) es denso.
        \item \(\varphi\) es continuo si, y sólo si, \(\forall c\in\R\) los conjuntos \(\{x:\varphi(x)<c\}\) y \(\{x:\varphi(x)>c\}\) son abiertos.
        \item Si \(A\subseteq E\) tiene interior no vacío y existe \(a\in\R\) tal que \(\varphi(x)\geq a\) para todo \(x\in A\), entonces \(\varphi\) es continuo.
    \end{enumerate}
\end{ejercicio}
\begin{proof}

\end{proof}

\vspace{0.5cm}

\begin{ejercicio}{3}
    Probar que los siguientes funcionales son lineales, continuos y hallar sus normas.
    \begin{enumerate}
        \item \(\varphi:c\to\C\), \(\varphi(x):=\displaystyle\lim_{n\to\infty}x_n\).
        \item \(\varphi:L^2([-1,1])\to\C\), \(\varphi(f):=\displaystyle\int_{-1}^1 tf(t)\,dt\).
        \item \(\varphi:\ell^\infty\to\C\), \(\varphi(x):=x_1+x_2\).
        \item \(\varphi:\ell^2\to\C\), \(\varphi(x):=x_1+x_2\).
    \end{enumerate}
\end{ejercicio}
\begin{proof}

\end{proof}

\vspace{0.5cm}

\begin{ejercicio}{4}
    Sea \(\varphi:(C[0,1],\norm{\cdot}_\infty)\to\R\) dada por
    \[
        \varphi(f)=\int_0^1 f(x)\,dx - f\left(\frac12\right).
    \]
    \begin{enumerate}
        \item Probar que \(\varphi\) es acotado y hallar su norma.
        \item Justificar, si existe o no, \(f\in C[0,1]\) tal que \(\norm{f}_\infty=1\) y \(|\varphi(f)|=2\).
    \end{enumerate}
\end{ejercicio}
\begin{proof}

\end{proof}

\vspace{0.5cm}

\begin{ejercicio}{5}
    Probar que la siguiente aplicación lineal \(T:\ell^p\to\ell^p\), \(1\leq p\leq\infty\), es continua y calcular su norma:
    \[
        T(x_1,x_2,x_3,\dots)=(x_n,x_{n-1},\dots,x_2,x_1,x_{n+1},x_{n+2},\dots).
    \]
\end{ejercicio}
\begin{proof}

\end{proof}

\vspace{0.5cm}

\begin{ejercicio}{6}
    Sea \(T:\ell^1\to\ell^1\) un operador lineal definido por
    \[
        T(x_1,x_2,x_3,\dots)=\left(2x_2,\frac32 x_3,\frac43 x_4,\dots\right).
    \]
    \begin{enumerate}
        \item Probar que \(T\) es acotado y hallar su norma.
        \item Más aún, hallar \(T^n(x_1,x_2,x_3,\dots)\) explícitamente y calcular \(\norm{T^n}\).
    \end{enumerate}
\end{ejercicio}
\begin{proof}

\end{proof}

\vspace{0.5cm}

\begin{ejercicio}{7}
    Sea \(\{a_n\}_{n\in\N}\) una sucesión acotada en un espacio de Banach \(E\). Demostrar que la aplicación \(T:\ell^1\to E\) definida por
    \[
        T(x)=\sum_{n\in\N} a_n x_n,
    \]
    es lineal, continua y \(\norm{T}=\sup_{n\in\N}\norm{a_n}\).
\end{ejercicio}
\begin{proof}

\end{proof}

\vspace{0.5cm}

\begin{ejercicio}{8}
    Sea \(E\) un espacio vectorial normado, probar que si \(\varphi,\psi\in E'\) son tales que \(\varphi\cdot\psi\equiv 0\), entonces \(\varphi\equiv 0\) o \(\psi\equiv 0\).
\end{ejercicio}
\begin{proof}

\end{proof}

\vspace{0.5cm}

\begin{ejercicio}{9}
    Operadores shift. Sea \(1\leq p<\infty\). Consideremos los operadores \(S,T:\ell^p\to\ell^p\) dados por
    \[
        S(x_1,x_2,x_3,\dots)=(0,x_1,x_2,x_3,\dots) \quad \text{(shift a derecha)}.
    \]
    \[
        T(x_1,x_2,x_3,\dots)=(x_2,x_3,x_4,\dots) \quad \text{(shift a izquierda)}.
    \]
    \begin{enumerate}
        \item Probar que \(S\in L(\ell^p)\) y es inyectivo. Calcular su norma.
        \item Probar que \(T\in L(\ell^p)\) y es suryectivo. Calcular su norma.
        \item Probar que \(TS=I\) pero \(ST\neq I\).
    \end{enumerate}
\end{ejercicio}
\begin{proof}

\end{proof}

\vspace{0.5cm}

\begin{ejercicio}{10}
    Operadores multiplicación.
    \begin{enumerate}
        \item Sea \(\varphi\in C([0,1])\). Consideremos el operador \(M_\varphi:C([0,1])\to C([0,1])\) dado por
              \[
                  M_\varphi(f)=\varphi f.
              \]
              Probar que \(M_\varphi\in L(C([0,1]))\) y calcular su norma.
        \item Sea \(\varphi\in L^\infty([0,1])\). Consideremos el operador \(M_\varphi:L^2([0,1])\to L^2([0,1])\) dado por
              \[
                  M_\varphi(f)=\varphi f.
              \]
              Probar que \(M_\varphi\in L(L^2([0,1]))\) y que
              \[
                  M_\varphi^2=M_\varphi \iff \varphi \text{ es una función característica.}
              \]
        \item Sea \(\{a_n\}_{n\in\N}\) una sucesión de números complejos. Si \(1\leq p<\infty\), definimos el operador \(A:\ell^p\to\ell^p\) dado por \(A(\{x_n\}_n)=\{x_n a_n\}_n\). Probar que:
              \begin{enumerate}
                  \item El operador \(A\) está bien definido si, y sólo si, \(\{a_n\}_{n\in\N}\in\ell^\infty\).
                  \item El operador \(A\) es inyectivo si, y sólo si, \(a_n\neq 0,\ \forall n\geq 1\).
                  \item \(A\) es un isomorfismo si, y sólo si, \(\{1/a_n\}_{n\in\N}\in\ell^\infty\).
              \end{enumerate}
    \end{enumerate}
\end{ejercicio}
\begin{proof}

\end{proof}

\vspace{0.5cm}

\begin{ejercicio}{11}
    Operadores integrales.
    \begin{enumerate}
        \item Sea \(E\) un espacio de Banach y \(T\in L(E)\). Probar que
              \begin{enumerate}
                  \item Si \(T\) está acotado inferiormente entonces \(R(T)\) es cerrado.
                  \item \(T\) está acotado inferiormente y es suryectivo si, y sólo si, \(T\) es inversible.
                  \item Si \(V:L^2[0,1]\to L^2[0,1]\) es el operador de Volterra dado por
                        \[
                            Vf(x)=\int_0^x f(t)\,dt,
                        \]
                        probar que \(V\) no está acotado inferiormente.
              \end{enumerate}
    \end{enumerate}
\end{ejercicio}
\begin{proof}

\end{proof}

\vspace{0.5cm}

\begin{ejercicio}{12}
    Sea \(E\) un espacio de Banach. Consideremos \(A,B\in L(E)\) y las sucesiones \(\{A_n\}_{n\in\N}\), \(\{B_n\}_{n\in\N}\in L(E)\). Probar que:
    \begin{enumerate}
        \item \(\norm{AB}\leq\norm{A}\norm{B}\).
        \item Si \(A_n\xrightarrow[n\to\infty]{} A\) y \(B_n\xrightarrow[n\to\infty]{} B\) entonces \(A_nB_n\xrightarrow[n\to\infty]{} AB\).
    \end{enumerate}
\end{ejercicio}
\begin{proof}

\end{proof}

\vspace{0.5cm}

\begin{ejercicio}{13}
    Sea \(E\) un espacio de Banach y sea \(A\in L(E)\) tal que \(\norm{A}<1\). Demostrar que \(I-A\) es inversible con
    \[
        (I-A)^{-1}=\sum_{n=0}^\infty A^n \quad \text{y} \quad \norm{(I-A)^{-1}}\leq\frac{1}{1-\norm{A}}.
    \]
\end{ejercicio}
\begin{proof}

\end{proof}

\vspace{0.5cm}

\begin{ejercicio}{14}
    Sea \(E\) un espacio de Banach. Probar que \(Gl(E):=\{T\in L(E):T \text{ es inversible}\}\) es abierto en \(L(E)\).

    \textit{Sugerencia}: sea \(A\in L(E)\) un operador inversible, probar que \(\bola{A}{\frac{1}{\norm{A^{-1}}}}\subset Gl(E)\).
\end{ejercicio}
\begin{proof}

\end{proof}

\vspace{0.5cm}

\begin{ejercicio}{15}
    Sea \(E\) un espacio de Banach. Consideremos una sucesión \(\{A_n\}_{n\in\N}\in L(E)\) de operadores inversibles tal que \(A_n\xrightarrow[n\to\infty]{} A\), para \(A\in L(E)\) no inversible. Probar que
    \[
        \norm{A_n^{-1}}\xrightarrow[n\to\infty]{}\infty.
    \]
\end{ejercicio}
\begin{proof}

\end{proof}

\end{document}
